\documentclass[11pt]{article}
\usepackage[T1]{fontenc}
\usepackage{lmodern}
\usepackage[margin=1in]{geometry}
\usepackage{amsmath,amssymb,amsthm,mathtools}
\usepackage{booktabs,longtable,array,microtype}
\usepackage{xcolor,xurl}
\usepackage[colorlinks=true,linkcolor=blue!50!black,citecolor=blue!50!black,urlcolor=blue!50!black]{hyperref}
\hypersetup{pdftitle={Smaller separating norms and a signed finalizer for Diophantine ZFC consistency},pdfauthor={Prepared for the ProveIt project},pdfsubject={Explicit Diophantine equations equivalent to ZFC consistency},pdfkeywords={Diophantine equations, ZFC consistency, Pell equations, Jones reduction, ProveIt}}
\usepackage{enumitem}
\setlist{itemsep=3pt,topsep=5pt}
\setlength{\emergencystretch}{2em}
\allowdisplaybreaks[2]
\newtheorem{theorem}{Theorem}[section]
\newtheorem{lemma}[theorem]{Lemma}
\newtheorem{proposition}[theorem]{Proposition}
\newtheorem{corollary}[theorem]{Corollary}
\theoremstyle{definition}\newtheorem{definition}[theorem]{Definition}
\theoremstyle{remark}\newtheorem{remark}[theorem]{Remark}
\newcommand{\NN}{\mathbb N}
\newcommand{\ZZ}{\mathbb Z}
\newcommand{\QQ}{\mathbb Q}
\newcommand{\PP}{\mathbb Z_{>0}}
\newcommand{\Con}{\operatorname{Con}}
\newcommand{\Norm}{\operatorname{Norm}}
\newcommand{\ZFC}{\mathrm{ZFC}}
\newcommand{\pin}{\texttt{ae3007ca4467011e0e051c84acb4a1dd13971b0e}}
\newcommand{\repo}[2]{\href{https://github.com/VladimirReshetnikov/ProveIt/blob/ae3007ca4467011e0e051c84acb4a1dd13971b0e/#1}{#2}}
\title{Smaller separating norms and a signed finalizer\newline for Diophantine ZFC consistency}
\author{Research synthesis prepared for the ProveIt project}
\date{23 October 2026}
\begin{document}
\maketitle
\begin{abstract}
We compare and combine the eight current Diophantine consistency packages in
ProveIt, numbered 16--23, at the fixed repository revision recorded below.
Two new algebraic reductions give concrete improvements. First, a polynomial
in twelve strictly positive integer unknowns has exact degree $1081$ and a
division-free circuit of $116$ binary arithmetic operations. Its
three-equation presentation has maximum degree $1072$ and $112$ operations.
This improves the former twelve-variable degree minimum $1145$ while matching
the former operation minimum $116$. The key change is the uniform bound
$|y'|<(T^2+XI^2)\sqrt G$ for every relevant algebraic conjugate, together with
a smaller coefficient for the first Pell root.

Second, a concrete Jones specialization in nine nonnegative unknowns has
exact degree $93264\,5^{33}+288$, compared with $161844\,5^{33}+484$ in
package 18, and uses $20\,547$ operations instead of $20\,531$. This is a
$42.374\%$ degree reduction at a cost of sixteen operations. It combines
structured radical bounds with a signed divisibility finalizer that avoids
squaring the divisor and dividend.

Both constructions use the complete, already published 943-tile coefficient
tuple of package 19, with $B=2^{5^{10}}$. We rebuild its four actual integers
and include their compressed binary representations. The nine-variable
index constants are defined by bounded finite recipes, rather than expanded.
Composing the new algebra with the repository's ordinary compiler and K215
calibration theorems proves in ZFC that each stated unsolvability assertion
is equivalent to $\Con(\ZFC)$. We give detailed proofs, exact polynomial
circuits, independent algebraic checks, and a reproducible archive. We claim
improvements within the inspected family, without claiming global minimality
or a new end-to-end proof-assistant certificate.
\end{abstract}

\noindent\textbf{Repository snapshot.} \pin. All comparisons use that snapshot;
the historical package numbers are retained for attribution.
\clearpage
\tableofcontents
\newpage
\input{sections/overview}
\input{sections/twelve}
\input{sections/nine}
\input{sections/coefficients}
\input{sections/verification}
\input{sections/references}
\end{document}
